% Compile with: xelatex practical_05_calc_02_extra.tex
% Armenian via polyglossia + Sylfaen (Windows). On Linux/Mac swap to "Noto Sans Armenian" or "DejaVu Sans".

\documentclass[11pt,a4paper]{article}

\usepackage{fontspec}
\setmainfont{Sylfaen}
\newfontfamily\armenianfont{Sylfaen}

\usepackage{polyglossia}
\setmainlanguage{armenian}
\setotherlanguage{english}

\usepackage{amsmath, amssymb, amsthm}
\usepackage[margin=2.5cm]{geometry}
\usepackage{hyperref}
\usepackage{tikz}
\usepackage{pgfplots}
\pgfplotsset{compat=1.18}
\usetikzlibrary{arrows.meta, calc, decorations.pathreplacing}

% Armenian flag palette
\definecolor{armred}{HTML}{D90012}
\definecolor{armblue}{HTML}{0033A0}
\definecolor{armorange}{HTML}{F2A800}

\newtheoremstyle{solstyle}
  {6pt}{6pt}{\normalfont}{}{\bfseries}{:}{0.5em}{}
\theoremstyle{solstyle}
\newtheorem*{solution}{Լուծում}
\newtheorem*{remark}{Դիտողություն}

\title{Թեյլորի բազմանդամ և $\dfrac{\sin x}{x}$ սահման}
\author{}
\date{}

\begin{document}
\maketitle

\section*{Խնդիր 1}
Կազմել $f(x) = \sin(x)$ ֆունկցիայի 4-րդ կարգի Թեյլորի բազմանդամը $a=0$ կետի լոկալ շրջակայքում:

\begin{solution}
$a$ կետի շրջակայքում $n$-րդ կարգի Թեյլորի բազմանդամի ընդհանուր բանաձևն է՝
\[
T_n(x) \;=\; \sum_{k=0}^{n} \frac{f^{(k)}(a)}{k!}\,(x-a)^k.
\]
Մեր դեպքում $a = 0$ (այսպիսի բազմանդամին ասում են նաև Մաքլորենի բազմանդամ), $n = 4$:

\medskip
Հաշվենք ածանցյալները 0 կետում.
\begin{align*}
f(x)       &= \sin(x),  & f(0)        &= 0, \\
f'(x)      &= \cos(x),  & f'(0)       &= 1, \\
f''(x)     &= -\sin(x), & f''(0)      &= 0, \\
f'''(x)    &= -\cos(x), & f'''(0)     &= -1, \\
f^{(4)}(x) &= \sin(x),  & f^{(4)}(0)  &= 0.
\end{align*}

\medskip
Տեղադրելով բանաձևի մեջ՝
\[
T_4(x) \;=\; 0 + 1\cdot x + \frac{0}{2!}x^2 + \frac{-1}{3!}x^3 + \frac{0}{4!}x^4
\;=\; x - \frac{x^3}{6}.
\]

\medskip
\textbf{Պատասխան՝} $\displaystyle T_4(x) = x - \frac{x^3}{6}$:
\end{solution}

\bigskip
\noindent\textbf{Տեսողական համեմատություն.}\quad Ստորև պատկերված է, թե ինչպես են $T_1(x)=x$ (գծային մոտարկում) և $T_3(x)=x-x^3/6$ (խորանարդային մոտարկում) ֆունկցիաները մոտարկում $\sin(x)$-ին 0-ի շրջակայքում: Որքան մոտ է $x$-ը 0-ին, այնքան ավելի լավ է մոտարկումը:

\begin{center}
\begin{tikzpicture}
\begin{axis}[
    width=12cm, height=7cm,
    axis lines=middle,
    xlabel={$x$}, ylabel={$y$},
    xmin=-4, xmax=4,
    ymin=-2.5, ymax=2.5,
    samples=200,
    legend pos=south east,
    legend cell align=left,
    legend style={font=\small, fill=white, fill opacity=0.85, draw=gray!50, text opacity=1},
    grid=both,
    grid style={gray!20},
    xtick={-3,-2,-1,1,2,3},
    ytick={-2,-1,1,2},
    every axis x label/.style={at={(ticklabel cs:1.02)}, anchor=west},
    every axis y label/.style={at={(ticklabel cs:1.02)}, anchor=south},
]
\addplot[armblue, very thick, domain=-3.5:3.5] {sin(deg(x))};
\addlegendentry{$\sin(x)$}
\addplot[armred, thick, dashed, domain=-3.5:3.5] {x};
\addlegendentry{$T_1(x) = x$}
\addplot[armorange, thick, domain=-3.5:3.5] {x - x^3/6};
\addlegendentry{$T_3(x) = x - x^3/6$}
\end{axis}
\end{tikzpicture}
\end{center}

\begin{remark}
Քանի որ $f^{(4)}(0) = 0$, ապա 4-րդ կարգի անդամը զրո է --- այսինքն, ստացված բազմանդամը համընկնում է նաև 3-րդ կարգի Թեյլորի բազմանդամի հետ: $\sin$-ի ամբողջական Մաքլորենի շարքն է՝
\[
\sin(x) \;=\; x - \frac{x^3}{3!} + \frac{x^5}{5!} - \frac{x^7}{7!} + \cdots
\;=\; \sum_{k=0}^{\infty} \frac{(-1)^k\, x^{2k+1}}{(2k+1)!}.
\]
\end{remark}

\bigskip
\section*{Խնդիր 2}
Ապացուցել, որ $\displaystyle \frac{\sin x}{x} \to 1$, երբ $x \to 0$:

\begin{solution}

\noindent\textbf{Մեթոդ 1 (Թեյլորի բանաձևով --- արագ մոտեցում).}\quad
Խնդիր 1-ից ունենք Թեյլորի վերլուծությունը մնացորդով՝
\[
\sin(x) \;=\; x - \frac{x^3}{6} + o(x^3), \qquad x \to 0,
\]
որտեղ $o(x^3)$ նշանակում է մի արտահայտություն, որի և $x^3$-ի հարաբերությունը ձգտում է զրոյի, երբ $x \to 0$:

\medskip
Բաժանելով $x$-ի ($x \neq 0$)՝
\[
\frac{\sin x}{x} \;=\; 1 - \frac{x^2}{6} + o(x^2).
\]

\medskip
Անցնելով սահմանի՝ $x \to 0$,
\[
\lim_{x \to 0} \frac{\sin x}{x}
\;=\; \lim_{x \to 0} \left(1 - \frac{x^2}{6} + o(x^2)\right)
\;=\; 1 - 0 + 0 \;=\; 1.
\]
\qed

\bigskip
\noindent\textbf{Մեթոդ 2 (Սեղմման թեորեմ --- դասական, ոչ շրջանաձև մոտեցում).}\quad
Թեյլորի մեթոդը պարունակում է թաքնված շրջանաձևություն. $\sin$-ի ածանցյալը 0 կետում --- այն, որ $\sin'(0)=1$ --- ինքը ապացուցվում է հենց $\lim_{x \to 0}\sin(x)/x = 1$ սահմանի միջոցով: Ուստի այս սահմանի \emph{հիմնարար} ապացույցը երկրաչափական է:

\medskip
Դիտարկենք $0 < x < \pi/2$ դեպքը: Միավոր շրջանագծի վրա վերցնենք $x$ ռադիան անկյունով սեկտոր, ինչպես պատկերված է ստորև:

\begin{center}
\begin{tikzpicture}[scale=4.2, >={Stealth[length=2mm]}]
  \def\ang{50}
  % Outer triangle (red filled, behind everything)
  \fill[armred!18] (0,0) -- (1,0) -- (1, {tan(\ang)}) -- cycle;
  % Sector (orange)
  \fill[armorange!35] (0,0) -- (1,0) arc[start angle=0, end angle=\ang, radius=1] -- cycle;
  % Inner triangle (blue, drawn on top of sector)
  \fill[armblue!40] (0,0) -- (1,0) -- ({cos(\ang)}, {sin(\ang)}) -- cycle;

  % Circle
  \draw[thick] (0,0) circle (1);
  % Outer triangle outline
  \draw[armred, very thick] (0,0) -- (1, {tan(\ang)}) -- (1,0);
  % Inner triangle outline
  \draw[armblue, very thick] (0,0) -- ({cos(\ang)}, {sin(\ang)}) -- (1,0);

  % Angle arc
  \draw[thick] (0.16, 0) arc[start angle=0, end angle=\ang, radius=0.16];
  \node at (0.27, 0.09) {\small $x$};

  % Vertices
  \node[below left]  at (0,0) {\small $O$};
  \node[below right] at (1,0) {\small $A$};
  \node[above right] at ({cos(\ang)}, {sin(\ang)}) {\small $B$};
  \node[right]       at (1, {tan(\ang)}) {\small $C$};

  % sin x bracket (vertical drop from B)
  \draw[armblue, thick, dashed] ({cos(\ang)}, {sin(\ang)}) -- ({cos(\ang)}, 0);
  \draw[<->, armblue, thick] ({cos(\ang)-0.05}, 0) -- ({cos(\ang)-0.05}, {sin(\ang)});
  \node[armblue, left] at ({cos(\ang)-0.05}, {sin(\ang)/2}) {\small $\sin x$};

  % tan x bracket
  \draw[<->, armred, thick] (1.1, 0) -- (1.1, {tan(\ang)});
  \node[armred, right] at (1.1, {tan(\ang)/2}) {\small $\tan x$};

  % Legend below
  \node[anchor=north west, font=\footnotesize] at (-0.3, -0.25) {
    \begin{tabular}{ll}
      \textcolor{armblue}{$\blacksquare$}  & ներքին եռանկյուն, մակերես $= \tfrac{1}{2}\sin x$ \\
      \textcolor{armorange}{$\blacksquare$} & շրջանի սեկտոր, մակերես $= \tfrac{1}{2}x$ \\
      \textcolor{armred}{$\blacksquare$}    & արտաքին եռանկյուն, մակերես $= \tfrac{1}{2}\tan x$
    \end{tabular}
  };
\end{tikzpicture}
\end{center}

\medskip
Համեմատելով՝ ներքին եռանկյան, շրջանի սեկտորի և արտաքին (շոշափողով կազմված) եռանկյան մակերեսները, ստանում ենք.
\[
\underbrace{\tfrac{1}{2}\sin x}_{\text{ներքին եռ.}}
\;\le\;
\underbrace{\tfrac{1}{2}\, x}_{\text{սեկտոր}}
\;\le\;
\underbrace{\tfrac{1}{2}\tan x}_{\text{արտաքին եռ.}}
\]
որից
\[
\sin x \;\le\; x \;\le\; \tan x \;=\; \frac{\sin x}{\cos x}.
\]

\medskip
Քանի որ $0 < x < \pi/2$ դեպքում $\sin x > 0$, բաժանենք $\sin x$-ի՝
\[
1 \;\le\; \frac{x}{\sin x} \;\le\; \frac{1}{\cos x},
\]
կամ համարժեքորեն, վերցնելով հակադարձ արժեքները (անհավասարությունները շրջվում են)՝
\[
\cos x \;\le\; \frac{\sin x}{x} \;\le\; 1.
\]

\medskip
$\sin$-ի կենտ և $x$-ի կենտ լինելու պատճառով $\dfrac{\sin x}{x}$ ֆունկցիան զույգ է, ուստի վերը նշված անհավասարությունը ճիշտ է նաև $-\pi/2 < x < 0$ դեպքում:

\medskip
Քանի որ $\cos$-ն անընդհատ է 0-ում և $\cos 0 = 1$, ապա $\displaystyle \lim_{x \to 0} \cos x = 1$: Կիրառելով սեղմման թեորեմը՝
\[
1 \;=\; \lim_{x \to 0} \cos x \;\le\; \lim_{x \to 0} \frac{\sin x}{x} \;\le\; \lim_{x \to 0} 1 \;=\; 1,
\]
հետևաբար
\[
\boxed{\;\lim_{x \to 0} \frac{\sin x}{x} \;=\; 1.\;}
\]
\qed
\end{solution}

\bigskip
\noindent\textbf{$\dfrac{\sin x}{x}$ ֆունկցիայի գրաֆիկը.}\quad Ֆունկցիան որոշված չէ $x=0$ կետում (բացված շրջանակը գրաֆիկում), բայց երբ $x$-ը մոտենում է 0-ին, ֆունկցիան անվերապահ ձգտում է 1-ին --- ինչը արտացոլվում է գրաֆիկի «շարունակական անցումով» սպիտակ կետի վրայով:

\begin{center}
\begin{tikzpicture}
\begin{axis}[
    width=12cm, height=6cm,
    axis lines=middle,
    xlabel={$x$}, ylabel={$y$},
    xmin=-8, xmax=8,
    ymin=-0.35, ymax=1.35,
    samples=400,
    grid=both,
    grid style={gray!20},
    xtick={-6.283, -3.14, 3.14, 6.283},
    xticklabels={$-2\pi$, $-\pi$, $\pi$, $2\pi$},
    ytick={-0.25, 0.5, 1},
    every axis x label/.style={at={(ticklabel cs:1.02)}, anchor=west},
    every axis y label/.style={at={(ticklabel cs:1.02)}, anchor=south},
]
\addplot[armred, thick, dashed, domain=-8:8, samples=2] {1};
\addplot[armblue, very thick, domain=-7.8:-0.04, samples=200] {sin(deg(x))/x};
\addplot[armblue, very thick, domain=0.04:7.8, samples=200] {sin(deg(x))/x};
\addplot[only marks, mark=*, mark options={fill=white, draw=armorange, line width=1pt}, mark size=2.5pt] coordinates {(0, 1)};
\node[armred, anchor=west, font=\small] at (axis cs:5.2, 1.12) {$y = 1$};
\node[anchor=south west, font=\small] at (axis cs:0.2, 1.02) {$\lim\limits_{x\to 0} = 1$};
\end{axis}
\end{tikzpicture}
\end{center}

\begin{remark}
Թվային ստուգում. $x = 0.1$ դեպքում՝ $\sin(0.1)/0.1 \approx 0{,}99833$, իսկ Թեյլորի 2-րդ կարգի մոտարկումից՝ $1 - 0{,}01/6 \approx 0{,}99833$: Համընկնում է մինչև 5 նիշ, ինչը հաստատում է $o(x^2)$ սխալի կարգը:
\end{remark}

\end{document}
